% --- Содержимое файла materials/6_2.tex ---

\subsection{Теорема Хана--Банаха}

\begin{idea}[Смысл теоремы]
	\textbf{Принцип сохранения нормы.}
	Представьте, что вы измеряете "длину" векторов только на плоскости внутри трехмерного пространства. У вас есть линейка (функционал).
	Теорема Хана-Банаха говорит, что вы можете использовать эту же линейку для измерения любых векторов в пространстве, при этом:
	1. Результаты измерений на плоскости не изменятся.
	2. "Масштаб" линейки (её норма) не увеличится.
	
	Это позволяет нам строить функционалы с нужными свойствами "по кирпичикам", начиная с малых подпространств.
\end{idea}

\begin{theorem}[Хана--Банаха]\label{Han-Banah}
	Пусть $E$ --- линейное нормированное пространство, $M \subset E$ --- линейное многообразие.
	Пусть $f$ --- линейный ограниченный функционал, заданный на $M$.
	
	Тогда существует функционал $\widetilde{f} \in E^*$ (определенный на всем $E$), который является продолжением $f$ с сохранением нормы:
	\begin{enumerate}
		\item $\widetilde{f}\big|_{M} = f$;
		\item $\|\widetilde{f}\|_{E^*} = \|f\|_{M^*}$.
	\end{enumerate}
\end{theorem}

\begin{proof}[Доказательство для вещественного сепарабельного случая]
	
	\textbf{Шаг 1. "Завоевание одного измерения".}
	Пусть $M \neq E$. Возьмем вектор $x_0 \notin M$. Рассмотрим подпространство
	$$
	M_1 = M \oplus [x_0].
	$$
	Любой элемент $y \in M_1$ однозначно представим в виде
	$$
	y = x + \alpha x_0,\quad x \in M.
	$$
	
	Зададим продолжение $f_1$ формулой:
	$$
	f_1(y) = f(x) + \alpha\, c,
	$$
	где
	$$
	c := f_1(x_0)
	$$
	— пока неизвестное число.
	
	Мы хотим выбрать $c$ так, чтобы норма не выросла:
	$$
	\|f_1\| \leqslant \|f\|
	$$
	(обратное неравенство очевидно). То есть для любого $y$:
	$$
	|f(x) + \alpha c| \leqslant \|f\| \cdot \|x + \alpha x_0\|.
	$$
	
	Если $\alpha = 0$, это верно. Если $\alpha \neq 0$, поделим на $|\alpha|$.
	Пусть $z = x/\alpha$. Неравенство примет вид:
	$$
	|f(z) + c| \leqslant \|f\| \cdot \|z + x_0\|.
	$$
	
	Раскроем модуль (двойное неравенство):
	$$
	-\|f\|\|z+x_0\| - f(z) \leqslant c \leqslant \|f\|\|z+x_0\| - f(z).
	$$
	
	Чтобы такое число $c$ существовало, нужно, чтобы \textbf{левая часть}
	(для любого $z_1$) была меньше или равна \textbf{правой части}
	(для любого $z_2$):
	$$
	-f(z_1) - \|f\|\|z_1+x_0\|
	\leqslant
	-f(z_2) + \|f\|\|z_2+x_0\|.
	$$
	
	Перегруппируем:
	$$
	f(z_2) - f(z_1)
	\leqslant
	\|f\|\bigl(\|z_1+x_0\| + \|z_2+x_0\|\bigr).
	$$
	
	Это действительно верно, так как:
	$$
	f(z_2) - f(z_1) = f(z_2 - z_1)
	\leqslant \|f\|\|z_2 - z_1\|
	= \|f\|\|(z_2+x_0) - (z_1+x_0)\|
	\leqslant \|f\|\bigl(\|z_2+x_0\| + \|z_1+x_0\|\bigr).
	$$
	
	Следовательно, существует число $c$, удовлетворяющее условию.
	Мы расширили функционал на одно измерение с сохранением нормы.
	
	\medskip
	\textbf{Шаг 2. Индукция (для сепарабельного случая).}
	
	Пусть $E$ сепарабельно. Выберем счетное всюду плотное множество $\{x_n\}$.
	
	Построим последовательность вложенных подпространств
	$$
	M_0 = M,\qquad
	M_n = M_{n-1} \oplus [x_{n-1}]
	$$
	(если $x_{n-1} \notin M_{n-1}$).
	
	На каждом шаге применяем Шаг 1 и получаем функционал $f_n$ с той же нормой.
	
	Определим функционал $f_\infty$ на объединении
	$$
	M_\infty = \bigcup M_n.
	$$
	Это линейное многообразие плотно в $E$.
	
	Так как функционал ограничен, по теореме о продолжении по непрерывности
	(Thm \ref{extrapolate-to-closure}) его можно единственным образом
	продолжить на всё $E$ с сохранением нормы.
\end{proof}

\subsubsection*{Следствия из Теоремы Хана-Банаха}

\begin{corollary}[О разделении точки и подпространства]\label{cor:HB_sep}
	Пусть $M \subsetneq E$ --- линейное подпространство, $x_0 \notin \overline{M}$.
	Тогда существует функционал $f \in E^*$, такой что:
	\begin{enumerate}
		\item $f(x) = 0$ для всех $x \in M$ ($M \subset \ker f$);
		\item $f(x_0) = 1$;
		\item $\|f\| = \frac{1}{\rho(x_0, M)}$.
	\end{enumerate}
\end{corollary}

\begin{idea}
	Это следствие говорит, что если точка не лежит в подпространстве (или его замыкании), то всегда найдется функционал, который "видит" эту разницу: он обнуляет всё подпространство, но на точке выдает единицу.
\end{idea}


\begin{proof}
	Рассмотрим $M_1 = M \oplus [x_0]$. Любой $y \in M_1$ имеет вид $y = x + \alpha x_0$.
	Зададим функционал $f_1$ на $M_1$: $f_1(x + \alpha x_0) = \alpha$.
	Очевидно, $f_1|_M = 0$ и $f_1(x_0) = 1$.
	
	Найдем норму $f_1$.
	\[ \|f_1\| = \sup_{y \neq 0} \frac{|f_1(y)|}{\|y\|} = \sup_{\alpha \neq 0, x \in M} \frac{|\alpha|}{\|x + \alpha x_0\|} = \sup \frac{1}{\|x/\alpha + x_0\|} = \frac{1}{\inf_{z \in M} \|z + x_0\|} = \frac{1}{\rho(x_0, M)}. \]
	По теореме Хана-Банаха продолжаем $f_1$ на всё $E$ с сохранением этой нормы.
\end{proof}

\begin{corollary}[Существование опорного функционала]\label{cor:HB_norm}
	Для любого $x \in E$ существует функционал $f \in E^*$, такой что:
	\begin{enumerate}
		\item $\|f\| = 1$;
		\item $f(x) = \|x\|$.
	\end{enumerate}
\end{corollary}

\begin{proof}
	Применим предыдущее следствие для $M = \{0\}$ и точки $x_0 = x$.
	Тогда $\rho(x, \{0\}) = \|x\|$.
	Существует $g \in E^*$ с нормой $1/\|x\|$, такой что $g(x) = 1$.
	Положим $f = \|x\| \cdot g$. Тогда $\|f\| = \|x\| \cdot (1/\|x\|) = 1$ и $f(x) = \|x\| \cdot 1 = \|x\|$.
\end{proof}

\begin{corollary}[Критерий равенства нулю]\label{hbcorollary3}
	Следующие утверждения эквивалентны:
	\begin{enumerate}
		\item $x = 0$;
		\item $\forall f \in E^*: f(x) = 0$.
	\end{enumerate}
\end{corollary}
\begin{proof}
	Если $x \neq 0$, то по следствию \ref{cor:HB_norm} существует $f$ с $f(x) = \|x\| \neq 0$. Противоречие.
\end{proof}

\begin{corollary}[Двойственность нормы]\label{hbcorollary4}
	Для любого $x \in E$:
	\[ \|x\| = \sup_{\|f\|_{E^*} \leqslant 1} |f(x)|. \]
\end{corollary}
\begin{proof}
	Неравенство $\leqslant$ очевидно: $|f(x)| \leqslant \|f\|\|x\| \leqslant \|x\|$.
	Равенство достигается на функционале из следствия \ref{cor:HB_norm} ($f(x)=\|x\|, \|f\|=1$).
\end{proof}